% !PRINTED 367
\weberpage{367}
%<<CONT
and since here $n = 11$, and thus odd, one must set:
\[
x = -5297, \quad y = -781.
\]
These are the absolutely smallest values of $x$, $y$ that satisfy equation (1).

The smallest positive solution is obtained from this by adding $\alpha$ and $\beta$ to them, thus
\[
x = 19038, \quad y = 2807.
\]
To this is also reduced the solution of the more general equation

\weq{4}{\alpha y - \beta x = \gamma}
in integers $x$, $y$.

First of all it is clear that, if (4) is to be solvable at all, every common divisor of $\alpha$ and $\beta$ must also be a divisor of $\gamma$. Such a common divisor is then removed by division. We therefore assume here too that $\alpha$ and $\beta$ have no common divisor. Then it follows, just as above, that all solutions of (4) are obtained from one of them, $x_0$, $y_0$, in the form
\[
x = x_0 + h \alpha, \quad y = y_0 + h \beta,
\]
in which $h$ is an undetermined integer. If one then sets in (4)
\[
x = \gamma \xi, \quad y = \gamma \eta,
\]
and divides by $\gamma$, then (4) passes into
\[
\alpha \eta - \beta \xi = 1,
\]
that is, into an equation of the form (1).

The solution becomes a completely determinate one when a further condition is given from which $h$ can be determined, e.g. the condition that $y$ shall lie between $0$ and $\beta$, with the inclusion of one of the two limits. We shall summarise this in the following theorem:

\gesperrt{III. The indeterminate equation}
\[
\alpha y - \beta x = \gamma
\]
\gesperrt{has, if $\alpha$, $\beta$, $\gamma$ are integers and $\alpha$ and $\beta$ are without common divisor, always an integral solution $x$, $y$, and only one if the condition}
\[
0 \leqq y < \beta,
\]
\gesperrt{or}
\[
0 < y \leqq \beta
\]
\gesperrt{is added.}
